\section{The Deployable Claim Is Transfer, and It Depends on What ``Unseen'' Means}
\label{sec:transfer}

If a learned model ties a parameter-free climatology on a city that has a radar archive, then the
only reason to prefer the model is that it runs where the climatology cannot be computed. A city
with no Sentinel-1 history has no climatology at all. So the question that matters is how well the
model does on ground it has never seen.

We previously reported held-out-city CSI of $0.094 \pm 0.021$ over three areas and called it a
floor rather than an estimate. With eleven further areas, being Bengaluru, the four remaining
Mumbai tiles and six Karnataka towns, adding 110 new Sentinel-1 masks, it was indeed a floor. The
reason it was so low is that a model trained on two areas has almost nothing to generalise from.

\subsection{The size of the correction depends entirely on how you hold data out}
Leaving out one \emph{tile} gives $0.270 \pm 0.015$, which is 12.5 times all-wet. That looks like a
threefold improvement, and it is misleading. With six Mumbai tiles and seven Karnataka areas in the
dataset, holding out Mumbai-West leaves five adjacent Mumbai tiles in training, and holding out
Chitradurga leaves six other Karnataka towns that share its climate and its tank morphology. That is
generalisation to unseen \emph{ground}, not to an unseen city.

So we group the areas into regions and hold out an entire region at a time. That gives
$0.131 \pm 0.005$, half the per-tile figure, and it is the number we report as the deployable one.
It is still 6.1 times all-wet and 1.4 times our earlier three-area floor, at a fifth of that
figure's seed spread.

\begin{table}[!t]
\centering
\caption{Transfer depends on what is held out. All terrain only, 3 seeds, identical training budget.
The rebuilt-stack control separates the effect of adding areas from the effect of rebuilding the
feature stack with the new builder's own definitions.}
\label{tab:transfer}
\begin{tabular}{lccc}
\toprule
held out & areas & mean CSI & $\times$ all-wet \\
\midrule
one tile (previously reported)  &  3 & 0.094 $\pm$ 0.021 & 3.6 \\
one tile, rebuilt stack control &  3 & 0.096 $\pm$ 0.009 & 3.7 \\
one tile                        & 14 & 0.270 $\pm$ 0.015 & 12.5 \\
one region                      & 14 & 0.131 $\pm$ 0.005 & 6.1 \\
\textbf{one region, with Chennai} & \textbf{15} & \textbf{0.147 $\pm$ 0.008} & \textbf{6.9} \\
\bottomrule
\end{tabular}
\end{table}

\subsection{The control that stops this being overstated}
Between the old result and the new one we also rewrote the dataset builder, so the feature stack
itself changed. That gives an alternative explanation for the gain which has nothing to do with
adding cities. We tested it directly by running the same three original areas on the rebuilt stack.
The answer is $0.096 \pm 0.009$ against the committed $0.094 \pm 0.021$. The rebuild moved nothing.
The gain is areas.

\subsection{Transfer is not uniform, and that is the operationally important part}
Held out entirely, the seven Karnataka areas score 0.213 at 10.3 times all-wet and Patna scores
0.166 at 3.1 times. The six Mumbai tiles score 0.033, which is 1.8 times all-wet and close to
predicting nothing at all.

Mumbai's Sentinel-1 water is tidal flat, creek and mangrove. Neither the Gangetic flood plain nor
the Karnataka plateau contains a single example of that kind of surface. That admits two readings
which fourteen areas cannot separate. Either coastal ground needs coastal training data, or tidal
water is not predictable from terrain at all. Section~\ref{sec:coastal} separates them.
